\documentclass[10pt]{beamer}

\usepackage[latin1]{inputenc}
\usepackage[spanish]{babel}

\mode<presentation>
\usetheme{Madrid}
%\usecolortheme[RGB={111,73,135}]{structure}
\usecolortheme[RGB={128,0,0}]{structure}
%\usecolortheme[RGB={0,96,0}]{structure}
%\usecolortheme[RGB={200,0,200}]{structure}
%\usecolortheme[RGB={0,128,0}]{structure}
%\usecolortheme[RGB={0,0,128}]{structure}
\usefonttheme{serif}
\useinnertheme{rectangles}
\useoutertheme{split}
\setbeamercovered{transparent}


% Definiciones para usar luego
\newtheorem{ejemplo}{Ejemplo}
\newtheorem{definicion}{Definici�n}

\title{Introduccion a AWK}
\author{Jes�s Espino Garc�a}
\date{10 de Octubre de 2008}
\subject{Introducci�n a AWK}

\institute[Kaleidos]{\includegraphics[height=1.5cm]{kaleidos}}

\begin{document}

  \frame{\maketitle}
  
  \section{Introducci�n}

  \begin{frame}{Un poco de historia}
    \begin{block}{Historia}
      \begin{itemize}
        \item Desarrollado por Alfred V. Aho, Peter J. Weinberger y Brian W. Kernighan
        \item Version original escrita en los laboratorios de AT\&T en 1977.
        \item En 1985 aparece una nueva version con muchas nuevas caracteristicas.
        \item En 1986 aparece la implementacion gnu de AWK, GAWK.
      \end{itemize}
    \end{block}
  \end{frame}
  
  \begin{frame}{�Qu� es AWK?}
    \begin{definicion}[AWK]
      AWK es equivalente a una navaja del ejercito suizo, que es �til para
      modificar archivos, buscar y transformar bases de datos, generar informes
      simples y otras muchas cosas.
    \end{definicion}
  \end{frame}

  \section{La entrada}

  \begin{frame}[containsverbatim]{Gesti�n de la entrada}
    \begin{block}{Tratamiento de la entrada}
      \begin{itemize}
        \item Cada fichero se divide en registros. Normalmente una registro por linea.
        \item Cada registro se divide en campos. Normalmente un campo por palabra.
        \item Los ficheros se procesan registro a registro.
        \item Para cada registro se generan las variables \verb+$0+ para todo el
              registro y \verb+$1+, \verb+$2+, \verb+$3+\dots para cada campos.
      \end{itemize}
    \end{block}
  \end{frame}
  
  \section{El proceso}

  \begin{frame}[containsverbatim]{Reglas en AWK}
    \begin{itemize}
      \item En AWK el c�digo se compone de reglas, divididas en:
	    \begin{block}{Patron}
		Un patron que se ajuste al registro. Ej: \verb+/hola/+.
	    \end{block}
	    \begin{block}{Accion}
		Un codigo a ejecutar. Ej: \verb+{print 'hola'}+.
	    \end{block}
      \item Si el patron es verdadero, se ejecuta la accion.
    \end{itemize}
  \end{frame}
  
  \begin{frame}[containsverbatim]{Patrones}
    \begin{block}{Patrones existentes}
      \begin{itemize}
        \item \verb+/expresion regular/+: Ejecuta la accion si el texto de
              entrada concuerda con la expresion regular. Ej: \verb+/^192.168/{print $0}+.
        \item \verb+expresion+: Ejecuta la acci�n si la expresi�n es verdadera.
              Ej: \verb+$1 == 192.168.0.1{print $0}+.
        \item \verb+pat1,pat2+: Ejecuta la acci�n para todas las cadenas entre la
              cadena que coincida con el primer patron y la que siguiente que coincida con el
              segundo patron. Ej: \verb+/^a/,/^b/{print $0}+.
        \item \verb+BEGIN+: Se ejecuta la acci�n antes de procesar ninguna
              entrada. Ej: \verb+BEGIN{print 'Principio de programa'}+.
        \item \verb+END+: Se ejecuta la acci�n despu�s de haber procesado toda la
              entrada. Ej: \verb+END{print 'Fin de programa'}+.
        \item \verb+null+: Ejecuta siempre la acci�n. Ej: \verb+{print $0}+
      \end{itemize}
    \end{block}
  \end{frame}
  
  \begin{frame}[containsverbatim]{Comparadores}
    \begin{block}{Comparadores}
      \begin{itemize}
        \item \verb+x > y+: Si x es mayor que y.
        \item \verb+x >= y+: Si x es mayor o igual que y.
        \item \verb+x < y+: Si x es menor que y.
        \item \verb+x <= y+: Si x es menor o igual que y.
        \item \verb+x == y+: Si x es igual a y.
        \item \verb+x != y+: Si x es diferente de y.
        \item \verb+x ~ y+: Si x encaja en la expresi�n regular y.
        \item \verb+x !~ y+: Si x no encaja en la expresi�n regular y.
      \end{itemize}
    \end{block}
  \end{frame}
  
  \begin{frame}[containsverbatim]{Operadores}
    \begin{block}{Operadores}
    \begin{itemize}
      \item \verb-x + y-: Suma.
      \item \verb+x - y+: Resta.
      \item \verb+- x+: Negaci�n.
      \item \verb+x * y+: Multiplicaci�n.
      \item \verb+x / y+: Divisi�n.
      \item \verb+x % y+: Resto.
      \item \verb+x ^ y+: Potencia.
      \item \verb+x ** y+: Potencia.
      \item \verb+x = y+: Asignacion.
    \end{itemize}
    \end{block}
  \end{frame}

  \begin{frame}[containsverbatim]{Sentencias de control}
    \begin{block}{Sentencias de control}
    \begin{itemize}
      \item \verb+if+: Bifurcacion if.
      \item \verb+while+: Bucle while.
      \item \verb+for+: Bucle for.
      \item \verb+break+: Sale del bucle actual.
      \item \verb+continue+: Salta al proximo ciclo del bucle.
      \item \verb+next+: Hace que pase a la siguiente cadena sin tener en cuenta el resto de las reglas.
      \item \verb+exit+: Hace que termine la ejecucion.
    \end{itemize}
    \end{block}
  \end{frame}
  
  \begin{frame}[containsverbatim]{Variables predefinidas para el sistema}
    \begin{block}{Variables predefinidas para el sistema}
    \begin{itemize}
      \item \verb+IGNORECASE+: Sensibilidad a may�sculas.
      \item \verb+FS+: Separador de campos.
      \item \verb+RS+: Separador de registros.
      \item \verb+OFS+: Separador de campos a la salida.
      \item \verb+ORS+: Separador de registros a la salida.
    \end{itemize}
    \end{block}
  \end{frame}
  
  \begin{frame}[containsverbatim]{Variables predefinidas para el usuario}
    \begin{block}{Variables predefinidas para el usuario}
    \begin{itemize}
      \item \verb+ARGC+: Numero de argumentos pasados.
      \item \verb+ARGV+: Argumentos pasados.
      \item \verb+ENVIRON+: Array con las variables de entorno.
      \item \verb+FILENAME+: El fichero que esta leyendo awk actualmente.
      \item \verb+FNR+: Numero del registro actual.
      \item \verb+NF+: Numero de campos del registro actual.
      \item \verb+NR+: Numero de registros procesados hasta el momento.
    \end{itemize}
    \end{block}
  \end{frame}
  
  \begin{frame}[containsverbatim]{Acciones}
    \begin{block}{Acciones}
    \begin{itemize}
      \item Cada acci�n es un c�digo de lo que debe hacer.
      \item Una de las principales funciones y mas usadas es \verb+print+.
      \item Existen muchas otras funciones predefinidas, podemos encontrar algunas en \verb+man awk+.
      \item Se pueden definir funciones propias.
    \end{itemize}
    \end{block}
  \end{frame}

  \section{Ejemplos}

  \begin{frame}[containsverbatim]{Ejemplos simples}
    \begin{ejemplo}[Codigos de una linea]
    \begin{itemize}
      \item \verb+{print $0}+ == \verb+cat+.
      \item \verb+/cadena/{print $0}+ == \verb+grep cadena+.
      \item \verb+(NR%2)==0{print $0}+: Imprime las lineas impares.
      \item \verb+BEGIN{FS=','}+: Establece que el separador de campo es ",".
      \item \verb+END{print 'Programa terminado'}+: Muestra un fin de programa al final
    \end{itemize}
    \end{ejemplo}
  \end{frame}

  \begin{frame}[containsverbatim]{Ejemplo complejo I}
    \begin{ejemplo}[Analisis de logs de apache]
    \begin{verbatim}
BEGIN{r200=0; r301=0; r304=0; r404=0;}
$9==200{r200++; next;}
$9==301{r301++; next;}
$9==304{r304++; next;}
$9==404{r404++; next;}
END{print "OK: " r200;
    print "Moved Permanently: " r301;
    print "Not Modified: " r304;
    print "Not Found: " r404;
    print "Others: " NR-r200-r301-r304-r404;
    print "Total: " NR}
    \end{verbatim}
    \end{ejemplo}
  \end{frame}
  
  \begin{frame}[containsverbatim]{Ejemplo complejo II}
    \begin{ejemplo}[Analisis de diccionario]
Ver el ejemplo dict.awk en la carpeta ejemplos.
    \end{ejemplo}
  \end{frame}
  
  \section{Para terminar.}

  \begin{frame}[containsverbatim]{Referencias}
    \begin{itemize}
      \item \verb+man awk+: Pagina del manual.
      \item http://www.gnu.org/software/gawk/manual/: Manual de AWK de GNU.
    \end{itemize}
  \end{frame}

  \begin{frame}{Dudas}
    \dots
  \end{frame}
  
  \begin{frame}
    \frametitle{\begin{center}Fin\end{center}}
    \begin{center}
      \includegraphics[height=3cm]{kaleidos}
    \end{center}
  \end{frame}

\end{document}
